\documentclass[11pt,a4paper]{article}

% Pakete
\usepackage[utf8]{inputenc}
\usepackage[T1]{fontenc}
\usepackage[ngerman]{babel}
\usepackage[left=18mm,right=18mm,top=18mm,bottom=18mm]{geometry}
\usepackage{helvet}
\renewcommand{\familydefault}{\sfdefault}
\usepackage{amsmath,amssymb}
\usepackage{graphicx}
\usepackage{tikz}
\usetikzlibrary{arrows.meta,positioning,patterns}
\usepackage{tcolorbox}
\tcbuselibrary{skins,breakable}
\usepackage{enumitem}
\usepackage{tabularx}
\usepackage{colortbl}
\usepackage{xcolor}
\usepackage[hidelinks]{hyperref}

% Fachfarbe Physik
\definecolor{physik}{HTML}{2c5aa0}
\definecolor{physiklight}{HTML}{e8f0fa}
\definecolor{loesung}{HTML}{c62828}

% Box-Stile
\tcbset{
    kasten/.style={
        colback=physiklight,
        colframe=physik,
        fonttitle=\bfseries,
        title=#1,
        boxrule=1pt,
        arc=2pt,
        left=5pt,
        right=5pt,
        top=5pt,
        bottom=5pt
    },
    formelbox/.style={
        colback=white,
        colframe=physik,
        boxrule=1pt,
        arc=2pt,
        left=6pt,
        right=6pt,
        top=6pt,
        bottom=6pt
    }
}

% Lücke mit Lösung
\newcommand{\lueckelsg}[1]{\textcolor{loesung}{\underline{\textbf{#1}}}}

% Kopfzeile
\usepackage{fancyhdr}
\pagestyle{fancy}
\fancyhf{}
\lhead{\small Physik 12 GK}
\chead{\small \textcolor{loesung}{MUSTERLÖSUNG}}
\rhead{\small Aufenthaltswahrscheinlichkeit}
\renewcommand{\headrulewidth}{0.4pt}

% Keine Einrückung
\setlength{\parindent}{0pt}
\setlength{\parskip}{6pt}

\begin{document}

% Titel
{\Large\textbf{Aufenthaltswahrscheinlichkeit}} \hfill \textcolor{loesung}{\textbf{MUSTERLÖSUNG}}\\[0.3em]
{\normalsize Fr 23.01.2026}

\vspace{0.3em}
\hrule height 1pt
\vspace{0.8em}

% ============================================================
% KASTEN 1: Stochastische Vorhersagbarkeit
% ============================================================

\begin{tcolorbox}[kasten={Kasten 1: Stochastische Vorhersagbarkeit}]

Ein einzelnes Elektron trifft am Doppelspalt an einem \lueckelsg{zufälligen} Ort auf.

Dieser Ort ist nicht \lueckelsg{vorhersagbar}.

\vspace{0.5em}

Schickt man jedoch \lueckelsg{viele} Elektronen durch den Doppelspalt,

entsteht ein \lueckelsg{Interferenzmuster}. Diese Verteilung ist vorhersagbar.

\vspace{0.8em}

\begin{center}
\fbox{\parbox{0.85\textwidth}{\centering\vspace{0.3em}
\textbf{Merksatz:} \glqq\lueckelsg{Einzeln zufällig}, in Summe \lueckelsg{determiniert}\grqq
\vspace{0.3em}}}
\end{center}

\end{tcolorbox}

\vspace{0.6em}

% ============================================================
% KASTEN 2: Wahrscheinlichkeitswelle
% ============================================================

\begin{tcolorbox}[kasten={Kasten 2: Die Wahrscheinlichkeitswelle}]

Die Wellenfunktion $\psi$ beschreibt nicht den genauen Ort eines Elektrons.

Sie beschreibt die \lueckelsg{Wahrscheinlichkeit}, das Elektron an einem bestimmten Ort zu finden.

\vspace{0.5em}

Die Größe $|\psi|^2$ heißt \lueckelsg{Aufenthaltswahrscheinlichkeitsdichte}.

\vspace{0.8em}

\textbf{Im Interferenzmuster:}
\begin{itemize}[leftmargin=1.5cm, itemsep=2pt]
\item Helle Streifen: $|\psi|^2$ ist \lueckelsg{groß} $\rightarrow$ hohe Trefferwahrscheinlichkeit
\item Dunkle Streifen: $|\psi|^2$ ist \lueckelsg{klein} $\rightarrow$ niedrige Trefferwahrscheinlichkeit
\end{itemize}

\end{tcolorbox}

\vspace{1em}

% ============================================================
% AUFGABE 1: Histogramm
% ============================================================

{\large\textbf{Aufgabe 1: Histogramm interpretieren}} \hfill \textbf{(5 BE)}

\vspace{0.5em}

\begin{enumerate}[label=\alph)}]
\item Beschreibe, wie sich das Histogramm verändert, wenn die Anzahl der Elektronen steigt. \hfill (2 BE)

\textcolor{loesung}{Bei wenigen Elektronen sieht das Histogramm unregelmäßig und zufällig aus. Mit steigender Anzahl wird das Interferenzmuster immer deutlicher erkennbar. Die Balken nähern sich der theoretischen $|\psi|^2$-Kurve an.}

\vspace{0.3em}

\item Erkläre, warum man bei 10 Elektronen noch kein klares Muster erkennt. \hfill (2 BE)

\textcolor{loesung}{Bei nur 10 Elektronen dominiert der Zufall. Jedes einzelne Elektron trifft an einem zufälligen Ort auf. Erst bei vielen Elektronen mittelt sich der Zufall heraus und die statistische Verteilung (das Interferenzmuster) wird sichtbar.}

\vspace{0.3em}

\item Was sagt $|\psi|^2$ über die Orte aus, an denen viele Elektronen auftreffen? \hfill (1 BE)

\textcolor{loesung}{An Orten, wo viele Elektronen auftreffen, ist $|\psi|^2$ groß. Das bedeutet: Die Aufenthaltswahrscheinlichkeit ist dort hoch.}

\end{enumerate}

\newpage

% ============================================================
% AUFGABE 2: Klassisch vs. Quantenmechanisch
% ============================================================

{\large\textbf{Aufgabe 2: Klassisch vs. Quantenmechanisch}} \hfill \textbf{(6 BE)}

\vspace{0.5em}

\begin{tabularx}{\textwidth}{|p{3.5cm}|X|X|}
\hline
\rowcolor{physiklight}
& \textbf{Kugeln (klassisch)} & \textbf{Elektronen (quantenmechanisch)} \\
\hline
\textbf{a) Erwartetes Muster auf dem Schirm}

\vspace{0.3em}
{\small (Skizze)}

\vspace{0.5cm}

&
\begin{center}
\textcolor{loesung}{
\begin{tikzpicture}[scale=0.6]
    \fill[loesung!50] (0,0) rectangle (0.4,2);
    \fill[loesung!50] (1,0) rectangle (1.4,2);
    \node[below] at (0.7,0) {\scriptsize 2 Streifen};
\end{tikzpicture}
}
\end{center}
&
\begin{center}
\textcolor{loesung}{
\begin{tikzpicture}[scale=0.6]
    \foreach \x/\h in {0/0.3, 0.3/0.8, 0.6/0.2, 0.9/1.5, 1.2/0.2, 1.5/0.8, 1.8/0.3} {
        \fill[loesung!50] (\x,0) rectangle (\x+0.2,\h);
    }
    \node[below] at (1,0) {\scriptsize Interferenz};
\end{tikzpicture}
}
\end{center}
\\
\hline
\end{tabularx}

\hfill {\small (je 2 BE)}

\vspace{0.8em}

\begin{enumerate}[label=\alph), start=2]
\item Erkläre den Unterschied zwischen den beiden Mustern. \hfill (2 BE)

\textcolor{loesung}{Bei Kugeln addieren sich die Wahrscheinlichkeiten einfach: Jede Kugel fliegt durch einen der beiden Spalte, das ergibt zwei Streifen hinter den Spalten.

Bei Elektronen interferieren die Wahrscheinlichkeitsamplituden (Wellenfunktionen). An manchen Stellen verstärken sie sich (helle Streifen), an anderen löschen sie sich aus (dunkle Streifen). Deshalb entstehen mehr als zwei Streifen.}

\end{enumerate}

\vspace{1em}

% ============================================================
% AUFGABE 3: Elektronenbeugung
% ============================================================

{\large\textbf{Aufgabe 3: Elektronenbeugung}} \hfill \textbf{(6 BE)}

\vspace{0.5em}

\begin{center}
\begin{tabular}{|c|c|c|c|c|}
\hline
\rowcolor{physiklight}
$U_B$ (kV) & $\lambda_{\text{theo}}$ (pm) & $r_1$ (cm) & $\lambda_{\text{exp}}$ (pm) & Abweichung (\%) \\
\hline
4,0 & \textcolor{loesung}{19,4} & \textcolor{loesung}{$\approx$ 1,23} & \textcolor{loesung}{$\approx$ 19} & \textcolor{loesung}{$< 5$} \\[0.4em]
\hline
6,0 & \textcolor{loesung}{15,8} & \textcolor{loesung}{$\approx$ 1,00} & \textcolor{loesung}{$\approx$ 16} & \textcolor{loesung}{$< 5$} \\[0.4em]
\hline
8,0 & \textcolor{loesung}{13,7} & \textcolor{loesung}{$\approx$ 0,87} & \textcolor{loesung}{$\approx$ 14} & \textcolor{loesung}{$< 5$} \\[0.4em]
\hline
\end{tabular}
\end{center}

\vspace{0.5em}

\begin{tcolorbox}[formelbox]
\textbf{Auswertungsformel (Kleinwinkelnäherung):} \quad $\lambda_{\text{exp}} \approx \dfrac{d \cdot r_1}{L}$

\vspace{0.3em}
{\small Gerätedaten: $d = 213$ pm (Netzebenenabstand), $L = 13{,}5$ cm (Abstand Folie–Schirm)}
\end{tcolorbox}

\vspace{0.3em}
{\small\textit{Hinweis: Die genauen Werte hängen von der Ablesegenauigkeit ab.}}

\vspace{0.8em}

\begin{enumerate}[label=\alph)}]
\item Wie verändert sich der Ringradius, wenn die Beschleunigungsspannung erhöht wird? \hfill (2 BE)

\textcolor{loesung}{Der Ringradius wird kleiner. Bei höherer Beschleunigungsspannung haben die Elektronen mehr Energie und einen größeren Impuls. Nach $\lambda = h/p$ wird die De-Broglie-Wellenlänge kleiner, und damit auch der Beugungswinkel und der Ringradius.}

\vspace{0.3em}

\item Erkläre in einem Satz, was dieses Experiment über Elektronen zeigt. \hfill (2 BE)

\textcolor{loesung}{Das Experiment zeigt, dass Elektronen Welleneigenschaften haben: Sie werden am Kristallgitter gebeugt, und ihre Wellenlänge entspricht der de Broglie-Formel $\lambda = h/p$.}

\end{enumerate}

\vspace{1em}

\begin{center}
\rule{10cm}{0.5pt}

\textbf{Gesamt:} \textcolor{loesung}{\textbf{17}} / 17 BE
\end{center}

\end{document}
